%%%%%%%%%%%%%%%%%%%%%%%%%%%%%%%%%
%
%       EXERCÍCIO
%
%%%%%%%%%%%%%%%%%%%%%%%%%%%%%%%%%

\definevalorquestao

\begin{exercicioBanco}[\valorquestao]
Um fazendeiro dispõe de material para construir 100 metros de cerca em uma região retangular, com um lado adjacente a um rio. Sabendo que ele não pretende colocar cerca no lado do retângulo adjacente ao rio, a área máxima da superfície que conseguirá cercar é de quanto?

% Define as alternativas
\newcommand{\alternativas}{%
    \begin{center}
        \begin{tabularx}{\textwidth}{XXXXX}
            (a) \(1000\) m\(^{2}\). &
            (b) \(1200\) m\(^{2}\). &
            (c) \(1250\) m\(^{2}\). &
            (d) \(1500\) m\(^{2}\). &
            (e) \(1600\) m\(^{2}\).
        \end{tabularx}
    \end{center}
}

% Define a resposta correta
\newcommand{\resposta}{C}

% Lógica condicional para exibição
\ifdefstring{\atividade}{lista}{%
    \alternativas
    \vspace{0.5em}
    
    \noindent\textbf{Resposta:} letra \textbf{\resposta}.
}{%
    \ifdefstring{\modo}{objetivo}{%
        \alternativas
    }{%
        % modo = discursiva → não mostra alternativas
    }
}
\end{exercicioBanco}

